\section{Loading data}\label{sec:loading}

This chapter covers both directions of the data flow: how to turn
quantum-chemistry outputs into a database, and how to read a database
without programming. It also defines the standard units and the information
each record holds, because every later chapter refers to them.

\subsection{Collecting program outputs into a database}\label{sec:collecting}

Build mode is used when every positional input is a quantum-chemistry output
file. The command is simply:

\begin{verbatim}
CJKQC <output files>... -o <database.jsonl>
\end{verbatim}

For each input, in the order given, \CJKQC{}:

\begin{enumerate}
  \item opens the file and looks at its leading lines to decide how to read
        it;
  \item reads the head and tail windows of the file;
  \item extracts the main calculated values of the final job (or final
        complete xyz frame);
  \item resolves the composition from the record's name (the file stem);
  \item stores exactly one record.
\end{enumerate}

The output database is replaced only after the whole run succeeds: either the
new database is installed, or an error leaves any previous database intact. A
failed or truncated job from a recognized format is still written, with status
\fld{failed} or \fld{truncated} and a short explanation in its error entry.

\begin{warning}
An input that matches no supported format, or an empty input, is not a job. It
aborts the run with exit code 1 and nothing is written. If you want to mix
databases and program outputs, merge them in a separate invocation.
\end{warning}

\subsubsection*{Windows and very large files}

\CJKQC{} does not read the whole of a large output file. It reads a head
window (default 256 KiB) and a tail window (default 4 MiB), which is where
banners, route sections, final energies, geometries, and frequencies live.
Files smaller than the two windows together are read whole. If a record's
error says that a table was cut by a window seam, widen the windows:

\begin{verbatim}
CJKQC huge.out --head-bytes 1048576 --tail-bytes 8388608 -o db.jsonl
\end{verbatim}

These flags affect build-mode reads only; merging databases ignores them.
A value below 1 is a usage error.

\subsubsection*{Reading several files at once}

By default collection reads one file at a time. On a shared cluster
filesystem, or when many files are being read, you can ask \CJKQC{} to work
on several files at once:

\begin{verbatim}
CJKQC *.log *.out --parallel 8 -o db.jsonl
\end{verbatim}

Each file is still handled as a whole, and the resulting database is identical
to the one produced by a one-at-a-time run, regardless of the number you give.
A value below 1 is a usage error. Merging databases always reads one file at a
time and ignores \flg{--parallel} for reading; if the run also uses
\flg{--arbalign}, that filtering step still uses the setting
(Section~\ref{sec:arbalign}).

\subsubsection*{Composition from file names}

The record's \fld{name} is the file stem, for example \texttt{2am3sa-1conf}
from \texttt{2am3sa-1conf.log}. When the stem follows the JKQC naming
convention, collection also stores the composition. Section
\ref{sec:names} explains the grammar. A name that does not follow the
convention is kept exactly as written and simply carries no composition
information.

\subsection{Merging databases}\label{sec:merging}

When every positional input is a database, \CJKQC{} runs in merge mode:

\begin{verbatim}
CJKQC part1.jsonl part2.jsonl -o combined.jsonl
\end{verbatim}

Records are concatenated in argument order and copied exactly as they were
stored: there is no re-reading, no composition resolution, and no rewriting of
the records. Information written by a newer version of \CJKQC{} therefore
survives a merge through an older version. Corrupt lines are skipped, reported
by path and line number, and make the exit code 1, but the valid records are
still written. \flg{--out} may name one of the inputs; all inputs are read
before the output is replaced.

\subsection{Reading a database file as a non-programmer}\label{sec:jsonl}

A database is a plain text file. Each line holds one complete record, written
as a set of names and values inside curly braces. The rules are deliberately
simple:

\begin{itemize}
  \item one record per line;
  \item every complete line ends with a line break;
  \item blank lines are ignored when the file is read;
  \item the \texttt{.jsonl} ending is only a convention; the content decides
        whether a file is a database;
  \item a corrupt or partially written line is skipped and reported, so the
        rest of the file remains usable;
  \item records may repeat; duplicate removal is a filtering decision, not a
        storage decision.
\end{itemize}

In the real file, each record is written on one long line. For readability,
the same record is shown below spread over several lines. Information that is
not available is left out completely; you will not see the word
\texttt{null}.

\begin{verbatim}
{
  "status": "ok",
  "name": "1dma",
  "source_path": "/scratch/runs/xtb/1dma.log",
  "program": "xtb",
  "program_version": "6.7.1",
  "command": "xtb 1dma.xyz --gxtb --acc 0.01 -ohess",
  "electronic_energy": -38.51,
  "charge": 0,
  "multiplicity": 1,
  "geometry": {
    "symbols": ["N", "H"],
    "coordinates": [0.0, 0.0, 0.0, 0.0, 0.0, 1.01]
  },
  "vibrational_frequencies": [3421.1],
  "components": {"dma": 1}
}
\end{verbatim}

What to notice:

\begin{itemize}
  \item \fld{status} says whether the job finished. If this entry is missing,
        it is treated as \fld{ok}.
  \item \fld{name} is the identity used for matching corrections and monomer
        references. \fld{source_path} is the full path of the original output
        file, when it could be determined.
  \item \fld{geometry.symbols} and \fld{geometry.coordinates} form two
        matching lists: one element name per atom, and three coordinates per
        atom in the same order. The coordinates are listed atom by atom: $x,y,z$
        for the first atom, then $x,y,z$ for the second, and so on.
  \item \fld{components} lists each component code together with its count.
        The codes are written alphabetically, so the same composition always
        looks the same.
  \item A \fld{correction} entry appears only after a correction has been
        attached. It never replaces \fld{electronic_energy}.
\end{itemize}

The entries may appear in any order, and unknown entries are ignored. New
versions of \CJKQC{} only add optional entries; they never rename or
repurpose an existing one. An older database therefore remains readable, and
a newer one remains readable by older tools.

\subsection{Record entries and standard units}\label{sec:fields}

The following table lists every entry a record can hold. Energies,
coordinates, frequencies, dipoles, and polarizabilities are always stored in
the standard units shown; conversion happens only when values are printed.

\begin{longtable}{@{}L{5.4cm}L{2.2cm}L{5.9cm}@{}}
\toprule
\textbf{Entry} & \textbf{Kind of value} & \textbf{Meaning and standard unit} \\
\midrule
\endhead
\fld{status} & text & \fld{ok}, \fld{failed}, or \fld{truncated}. \\
\fld{error} & text & Plain-language reason for a failed or truncated job;
present whenever the status is \fld{failed}. \\
\fld{name} & text & Structure identity: the source file's stem, exactly as
written. \\
\fld{source_path} & text & Full path of the original output file. \\
\fld{program} & text & Program name: \fld{gaussian}, \fld{orca}, \fld{xtb},
\fld{mrcc}, or \fld{xyz}. \\
\fld{program_version} & text & Version string as printed by the program. \\
\fld{command} & text & The run's command line as echoed by the program. \\
\fld{electronic_energy} & number & Total electronic energy of the final
structure, in hartree. \\
\fld{scf_energy} & number & Converged SCF energy of the final structure, in
hartree. Used only by extrapolation. \\
\fld{charge} & whole number & Net charge, in elementary charges. \\
\fld{multiplicity} & whole number & Spin multiplicity, at least 1. \\
\fld{geometry} & two lists & \fld{symbols} (element names) and
\fld{coordinates} (angstrom, listed atom by atom). \\
\fld{vibrational_frequencies} & list & Genuine vibrational modes in
cm$^{-1}$; a negative value marks an imaginary mode. \\
\fld{gradient} & list & Cartesian gradient in hartree/\AA{}, listed atom by
atom, three values per atom. \\
\fld{dipole_moment} & list & Dipole vector $[x,y,z]$ in atomic units
($e\,a_0$). \\
\fld{polarizability} & list & The nine polarizability tensor values in
atomic units, in the order $xx, xy, xz, yx, yy, yz, zx, zy, zz$. The
isotropic value is calculated from them when printed. \\
\fld{mulliken_charges} & list & Per-atom Mulliken charges in elementary
charges, in geometry atom order. \\
\fld{rotational_symmetry_number} & whole number & Point-group rotational
symmetry number $\sigma$, at least 1. \\
\fld{components} & code and count & Resolved composition. \\
\fld{correction} & group of values & Attached single-point correction and its
origin; see Chapter~\ref{sec:corrections}. \\
\bottomrule
\end{longtable}

The standard units are summarized here once:

\begin{center}
\begin{tabular}{@{}ll@{}}
\toprule
Quantity & Unit stored in the database \\
\midrule
Energy, enthalpy, Gibbs, zero-point energy & hartree \\
Coordinates and distances & \AA{} (angstrom) \\
Gradient & hartree/\AA{} \\
Vibrational frequencies & cm$^{-1}$ \\
Dipole moment & atomic units, $e\,a_0$ \\
Polarizability & atomic units, $e^2 a_0^2/E_\mathrm{h}$ \\
Charge and Mulliken charges & elementary charge $e$ \\
Mass & unified atomic mass unit, u \\
Multiplicity and rotational symmetry number & whole numbers \\
\bottomrule
\end{tabular}
\end{center}

Printing in kcal/mol is a presentation choice made with \flg{--unit}; it
never changes the database.

\subsection{Names and compositions}\label{sec:names}

The JKQC naming convention encodes a composition in the file stem:

\begin{center}
\begin{tabular}{@{}ll@{}}
\toprule
Name & Resolved components \\
\midrule
\texttt{1dma} & \texttt{\{"dma": 1\}} \\
\texttt{2cha} & \texttt{\{"cha": 2\}} \\
\texttt{1pinic1sa} & \texttt{\{"pinic": 1, "sa": 1\}} \\
\texttt{2am3sa-1conf} & \texttt{\{"am": 2, "sa": 3\}} \\
\texttt{1am1p} & \texttt{\{"am": 1, "p": 1\}} \\
\texttt{h2o} & no composition (does not match the grammar) \\
\texttt{movie} & no composition \\
\bottomrule
\end{tabular}
\end{center}

The grammar is:

\begin{verbatim}
name        := composition ( "-" label )?
composition := ( count code )+        (at least one pair)
count       := digit+                  (1 to 4294967295)
code        := letter+                 (plain letters, either case)
label       := any text                (never read, never stored)
\end{verbatim}

Details worth remembering:

\begin{itemize}
  \item The name splits at its \emph{first} hyphen. Everything after it is a
        free-form label and is ignored.
  \item Counts allow leading zeros: \texttt{02am} means two \texttt{am}.
  \item Codes are case-sensitive: \texttt{am} and \texttt{Am} are different
        components.
  \item Repeated codes sum: \texttt{1am1am} resolves to
        \texttt{\{"am": 2\}}.
  \item A count of zero, a special character, a digit inside a code, or an
        implicit count such as \texttt{am2sa} makes the name unresolvable.
        This is never an error; the record simply has no \fld{components}.
\end{itemize}

Later features use the stored composition, not the raw name. This is
why a database written once can be filtered and post-processed without
re-reading file names.

\subsection{Record statuses}\label{sec:statuses}

\begin{description}
  \item[\fld{ok}] The job converged and its values are final.
  \item[\fld{failed}] The job did not finish. Some values may still be
        present, and \fld{error} explains what happened.
  \item[\fld{truncated}] The output file ends in the middle of the run. Some
        values may still be present.
\end{description}

The filtering steps drop all non-\fld{ok} records before anything else and
report their count as a warning. A print-only run prints the
records as loaded, so a failed record can still appear in a table unless you
filter it out.
